\documentclass[11pt,a4paper,sans]{moderncv}

% Estilo moderncv
\moderncvstyle{classic} % opciones: 'casual', 'classic', 'banking', 'oldstyle', 'fancy'
\moderncvcolor{blue} % opciones: 'black', 'blue', 'burgundy', 'green', 'grey', 'orange', 'purple', 'red'

% Paquetes
\usepackage[utf8]{inputenc}
\usepackage[T1]{fontenc}
\usepackage[spanish]{babel}
\usepackage[scale=0.85]{geometry}
\usepackage{lmodern} % Fuente escalable

% Datos personales
\name{Fabián}{Belmar}
\title{Doctor en Procesos e Instituciones Políticas}
\email{fbelmar@cepchile.cl}
\phone[mobile]{+56 9 76371608}
\homepage{belmarfabian.github.io}

\begin{document}

\makecvtitle

% ============================================
% FORMACIÓN ACADÉMICA
% ============================================
\section{Formación Académica}

\cventry{2020--2025}{Doctorado en Procesos e Instituciones Políticas}{Universidad Adolfo Ibáñez}{Chile}{Distinción máxima}{Tesis: ``Corrupción al margen del Estado. El caso de la FIFA''. Director: Dr. Aldo Mascareño}

\cventry{2013--2018}{Licenciado en Ciencia Política y Administración Pública}{Universidad de Talca}{Chile}{Distinción máxima}{Tesis: ``Vinculación no programática en la gestión municipal''. Director: Dr. Mauricio Morales}

% ============================================
% EXPERIENCIA ACADÉMICA
% ============================================
\section{Experiencia Académica}

\cventry{2025--}{Investigador Asistente}{Centro de Estudios Públicos (CEP)}{Santiago}{}{Proyecto C22 - Métodos Digitales. Análisis de procesos sociales mediante ciencia de datos y machine learning.}

\cventry{2023--2024}{Investigador Pasante}{Instituto Gino Germani, U. Buenos Aires}{Argentina}{}{Grupo de Estudios de Culturas Populares Contemporáneas. Directora: Dra. Verónica Moreira.}

\cventry{2022--2023}{Asistente de Investigación}{FONDECYT Regular \#1220004}{}{}{Proyecto ``Vinculación entre representantes y representados''. IP: Dr. Mauricio Morales.}

\cventry{2018--2023}{Coordinador de Proyectos}{Centro de Análisis Político, U. de Talca}{}{}{Investigación, vinculación con el medio, análisis de datos.}

% ============================================
% PUBLICACIONES
% ============================================
\section{Artículos Científicos (Selección)}

\cvitem{2025}{\textbf{Belmar, F.}, \& Mascareño, A. Corruption: An uneven field of research. \textit{Societies}, 15(7). [WoS/Scopus]}

\cvitem{2025}{\textbf{Belmar, F.}, \& Morales, M. Ethnic solidarity in the non-programmatic distribution of benefits. \textit{Ethnic and Racial Studies}. [WoS/Scopus]}

\cvitem{2024}{\textbf{Belmar, F.}, Lara, B., \& Miguieles, J. Unpacking electoral competition. \textit{Policy Studies}. [WoS/Scopus]}

\cvitem{2024}{\textbf{Belmar, F.}, Maldonado, F., \& Morales, M. In search of the missing link. \textit{Public Integrity}. [WoS/Scopus]}

\cvitem{2023}{\textbf{Belmar, F.}, Morales, M., \& Villarroel, B. Writing constitution without parties? \textit{Politics}. [WoS/Scopus]}

\cvitem{2023}{\textbf{Belmar, F.}, et al. Demanding non-programmatic distribution. \textit{Policy Studies}. [WoS/Scopus]}

\cvitem{2022}{\textbf{Belmar, F.}, \& Morales, M. Clientelism, turnout and incumbents' performance. \textit{Social Sciences}. [WoS/Scopus]}

\cvitem{2020}{\textbf{Belmar, F.}, \& Morales, M. Clientelismo en la gestión local. \textit{Política y Sociedad}. [WoS/Scopus]}

% ============================================
% CAPÍTULOS DE LIBRO
% ============================================
\section{Capítulos de Libro}

\cvitem{2024}{Descentralización y participación ciudadana en Chile. En \textit{La Revolución Silente}.}

\cvitem{2020}{Las tecnologías al servicio de la comunicación parlamentaria. En \textit{Parlamentos Conectados}.}

\cvitem{2018}{El efecto de la elección de concejales sobre consejeros regionales. En \textit{Política Subnacional en Chile}.}

% ============================================
% DOCENCIA
% ============================================
\section{Docencia}

\cvitem{2025}{U. de Talca: Evaluación de Políticas; Ciencia de la Administración Pública.}
\cvitem{2025}{U. Diego Portales: Métodos Cuantitativos Introductorios y Avanzados.}
\cvitem{2023--2024}{U. Central: Análisis de Datos Cuantitativos; Métodos Cuantitativos.}
\cvitem{2021--2022}{U. de Talca: Métodos Cualitativos; Análisis de Datos.}

% ============================================
% BECAS Y DISTINCIONES
% ============================================
\section{Becas y Distinciones}

\cvitem{2023--2024}{Beca Pasantía Internacional ANID (12,000 USD)}
\cvitem{2021--2022}{Beca Beneficios Complementarios ANID (6,000 USD)}
\cvitem{2020}{Beca Doctorado Nacional ANID (8,000 USD)}
\cvitem{2019}{Fondo de Desarrollo Institucional, Ministerio de Educación (7,000 USD)}
\cvitem{2018}{Mejor Egresado de la Generación, Universidad de Talca}

% ============================================
% HABILIDADES
% ============================================
\section{Habilidades}

\cvitem{Programación}{R (Avanzado), Python (Intermedio), SQL (Intermedio), \LaTeX\ (Avanzado)}
\cvitem{Análisis}{SPSS, STATA, Jamovi, QGIS, ATLAS.ti (Avanzado)}
\cvitem{Idiomas}{Español (Nativo), Inglés (Intermedio escrito)}

% ============================================
% REFERENCIAS
% ============================================
\section{Referencias}

\cvitem{}{Dr. Mauricio Morales, Universidad de Talca. \emailsymbol mmoralesq@utalca.cl}
\cvitem{}{Dr. Aldo Mascareño, Universidad Adolfo Ibáñez. \emailsymbol amascareno@cepchile.cl}

\end{document}
